\chapter{Le variabili aleatorie continue notevoli}

Come nel discreto avevamo la binomiale, l'uniforme e la Poisson, anche nel continuo ci sono alcune variabili che
compaiono di continuo. In questo capitolo ne studiamo quattro: l'\textbf{uniforme}, l'\textbf{esponenziale}, quella
di \textbf{Laplace} e quella di \textbf{Cauchy}. Per ognuna calcoliamo la densità, la CDF e la media; per farlo ci
servono alcuni strumenti nuovi per gli integrali (l'integrazione per parti, il cambio di variabile) e due funzioni
nuove (il logaritmo e l'arcotangente), tutti spiegati da zero. La variabile di Cauchy ci riserverà una sorpresa: non ha
media. La quinta variabile continua notevole, la gaussiana, è così importante che avrà un capitolo tutto per sé.

% ============================================================
\section{Due funzioni utili: il gradino e il rettangolo}
% ============================================================

Le densità che vedremo sono spesso diverse da zero solo su una parte della retta. Per scriverle in una sola formula
si usano due funzioni «interruttore».

\begin{richiamo}{il gradino unitario e la funzione rettangolo}
Il \term{gradino unitario} $u(x)$ vale 0 per $x < 0$ e 1 per $x \geq 0$: moltiplicare una funzione per $u(x)$ la
«spegne» a sinistra dello zero. La \term{funzione rettangolo} $\mathrm{rect}(x)$ vale 1 per $\abs{x} < \frac12$ e 0
altrove: è un impulso largo 1 centrato in 0, e $\mathrm{rect}\left(x - \frac12\right)$ è lo stesso impulso spostato
sull'intervallo $(0, 1)$.
\begin{center}
\begin{tikzpicture}
\begin{groupplot}[group style={group size=3 by 1, horizontal sep=1cm},
  quaderno, width=5cm, height=3.4cm, xmin=-1.6, xmax=1.8, ymin=0, ymax=1.45, ytick={1}, xtick={-1,0,1},
  title style={font=\small}]
\nextgroupplot[title={$u(x)$}, xtick={0}]
  \addplot[mspiric, line width=1.3pt] coordinates {(-1.5,0) (0,0)};
  \addplot[mspiric, line width=1.3pt] coordinates {(0,1) (1.6,1)};
  \addplot[only marks, mark=*, mark size=1.6pt, mspiric] coordinates {(0,1)};
\nextgroupplot[title={$\mathrm{rect}(x)$}, xtick={-0.5,0.5}, xticklabels={$-\frac12$,$\frac12$}]
  \addplot[mspiric, line width=1.3pt] coordinates {(-1.5,0) (-0.5,0) (-0.5,1) (0.5,1) (0.5,0) (1.6,0)};
\nextgroupplot[title={$\mathrm{rect}\left(x - \frac12\right)$}, xtick={0,1}]
  \addplot[mspiric, line width=1.3pt] coordinates {(-1.5,0) (0,0) (0,1) (1,1) (1,0) (1.6,0)};
\end{groupplot}
\end{tikzpicture}
\end{center}
Per esempio, $e^{-x} u(x)$ vale $e^{-x}$ per $x \geq 0$ e zero per $x < 0$. Quanto valga la funzione esattamente nel
punto di salto non conta: un singolo punto non ha area, quindi non cambia nessun integrale.
\end{richiamo}

% ============================================================
\section{La variabile uniforme}
% ============================================================

\dfn{Variabile uniforme continua}{
  Una variabile aleatoria $X$ si dice \term{uniformemente distribuita} nell'intervallo $[a, b]$ (con $b > a$), in breve
  $X \sim \mathcal{U}(a, b)$, se
  \[ f_X(x) = \begin{cases} \dfrac{1}{b - a} & x \in (a, b) \\[2mm] 0 & \text{altrove} \end{cases} \]
}

La densità è costante sull'intervallo: nessun punto è privilegiato. L'altezza $\frac{1}{b-a}$ è l'unica che rende
l'area del rettangolo (base $b - a$) uguale a 1. Il supporto è $\mathrm{supp}[f_X(x)] = [a, b]$, e tale è anche
l'alfabeto: $X$ non assume valori fuori dall'intervallo. Con la funzione rettangolo,
$\mathcal{U}(0, 1)$ ha densità $\mathrm{rect}\left(x - \frac12\right)$.

\teorema[uniformec]{CDF e media dell'uniforme}{
  Se $X \sim \mathcal{U}(a, b)$:
  \[ F_X(x) = \begin{cases} 0 & x < a \\[1mm] \dfrac{x - a}{b - a} & a \leq x \leq b \\[2mm] 1 & x > b
     \end{cases} \qquad\qquad E[X] = \frac{a + b}{2} . \]
}

\begin{dimostrazione}
\passo{La CDF} A sinistra di $a$ non c'è area. Fra $a$ e $x$ la densità vale $\frac{1}{b-a}$, costante, quindi l'area
è quella di un rettangolo: base $x - a$, altezza $\frac{1}{b-a}$. Oltre $b$ l'area è tutta, cioè 1.
\[ F_X(x) = \int_{-\infty}^{x} f_X(t) \dd t = \int_a^x \frac{1}{b - a} \dd t = \frac{x - a}{b - a} \qquad (a \leq x \leq b) . \]
\passo{La media} Con la primitiva di $x$, che è $\frac{x^2}{2}$:
\[ E[X] = \int_a^b x\, \frac{1}{b - a} \dd x = \frac{1}{b - a} \left[ \frac{x^2}{2} \right]_a^b = \frac{b^2 - a^2}{2(b
   - a)} = \frac{(b - a)(b + a)}{2(b - a)} = \frac{b + a}{2} , \]
dove si è usato il prodotto notevole $b^2 - a^2 = (b - a)(b + a)$.
\end{dimostrazione}

\begin{figure}[H]
\centering
\begin{tikzpicture}
\begin{groupplot}[group style={group size=2 by 1, horizontal sep=1.8cm},
  quaderno, width=6.6cm, height=4.2cm, xmin=-0.6, xmax=4.4, ymin=0, ymax=1.25, xtick={1,2,3}, xticklabels={$a$,$\frac{a+b}{2}$,$b$}]
\nextgroupplot[ylabel=$f_X(x)$, xlabel=$x$, ytick={0.5}, yticklabels={$\frac{1}{b-a}$}]
  \addplot[fill=mspiA!20, draw=none] coordinates {(1,0) (1,0.5) (3,0.5) (3,0)} \closedcycle;
  \addplot[mspiA, line width=1.2pt] coordinates {(-0.4,0) (1,0) (1,0.5) (3,0.5) (3,0) (4.2,0)};
  \draw[mspiB, dashed] (axis cs:2,0) -- (axis cs:2,0.75) node[above, font=\scriptsize] {media};
\nextgroupplot[ylabel=$F_X(x)$, xlabel=$x$, ytick={0.5,1}, yticklabels={$\frac12$,$1$}]
  \addplot[mspiB, line width=1.2pt] coordinates {(-0.4,0) (1,0) (3,1) (4.2,1)};
\end{groupplot}
\end{tikzpicture}
\caption{La variabile uniforme: densità costante sull'intervallo, CDF che sale in linea retta da 0 a 1.}
\end{figure}

% ============================================================
\section{La variabile esponenziale}
% ============================================================

\dfn{Variabile esponenziale}{
  Una variabile aleatoria $X$ si dice \term{esponenziale} (unilatera) di parametro $\lambda > 0$, in breve
  $X \sim \mathcal{E}(\lambda)$, se ha densità
  \[ f_X(x) = \lambda\, e^{-\lambda x}\, u(x) . \]
}

Il supporto è $\mathrm{supp}[f_X(x)] = [0, \infty)$: la variabile è sempre non negativa, $X \geq 0$. È la variabile
tipica dei \textbf{tempi di attesa}: il tempo fino al prossimo pacchetto che arriva a un router, fino al guasto di un
componente, fino al decadimento di un atomo. La densità è massima in 0 e cala esponenzialmente: le attese brevi sono le
più probabili, quelle lunghe sono rare. Prima di calcolare la CDF serve sapere come si integra un esponenziale.

\begin{richiamo}{la derivata e la primitiva di $e^{ax}$}
Per una costante $a$: $\left(e^{ax}\right)' = a\, e^{ax}$. Infatti il rapporto incrementale è
\[ \frac{e^{a(x+h)} - e^{ax}}{h} = e^{ax}\, \frac{e^{ah} - 1}{h} = a\, e^{ax}\, \frac{e^{ah} - 1}{ah} , \]
e l'ultima frazione tende a 1 quando $h \to 0$ (è il limite $\frac{e^k - 1}{k} \to 1$ del Capitolo~7, con $k = ah$ che
tende anch'esso a zero). Di conseguenza una primitiva di $e^{ax}$ è $\frac{e^{ax}}{a}$, e in particolare
\[ \int e^{-\lambda t} \dd t \quad\text{ha primitiva}\quad \frac{e^{-\lambda t}}{-\lambda} . \]
Serve anche sapere che l'esponenziale «batte» qualsiasi potenza: per $\lambda > 0$, $x^n e^{-\lambda x} \to 0$ quando
$x \to \infty$. Per $n = 1$: dalla serie esponenziale, che ha tutti i termini positivi, $e^{\lambda x} \geq
\frac{(\lambda x)^2}{2}$, quindi $x\, e^{-\lambda x} \leq \frac{2x}{\lambda^2 x^2} = \frac{2}{\lambda^2 x} \to 0$. Per
$n$ qualsiasi si ragiona allo stesso modo con il termine $\frac{(\lambda x)^{n+1}}{(n+1)!}$ della serie.
\end{richiamo}

\teorema[esponenzialecdf]{CDF dell'esponenziale}{
  Se $X \sim \mathcal{E}(\lambda)$: \qquad $F_X(x) = \left( 1 - e^{-\lambda x} \right) u(x)$, \qquad
  $\overline{F}_X(x) = P(X > x) = e^{-\lambda x}$ per $x \geq 0$.
}

\begin{dimostrazione}
\passo{} Per $x < 0$ non c'è area. Per $x \geq 0$, con la primitiva del richiamo:
\[ F_X(x) = \int_0^x \lambda\, e^{-\lambda t} \dd t = \lambda \left[ \frac{e^{-\lambda t}}{-\lambda} \right]_0^x =
   \left[ -e^{-\lambda t} \right]_0^x = -e^{-\lambda x} - (-e^0) = 1 - e^{-\lambda x} . \]
La CCDF è il complemento a 1: $e^{-\lambda x}$.
\end{dimostrazione}

Per la media serve una tecnica nuova, perché l'integrando $x\, \lambda e^{-\lambda x}$ è un prodotto.

\begin{richiamo}{l'integrazione per parti}
La regola del prodotto (Capitolo~7) dice che $(f g)' = f' g + f g'$. Integrando entrambi i membri fra $a$ e $b$, e
ricordando che l'integrale di una derivata è la differenza dei valori agli estremi (teorema fondamentale):
\[ \Big[ f g \Big]_a^b = \int_a^b f'(x)\, g(x) \dd x + \int_a^b f(x)\, g'(x) \dd x . \]
Portando un integrale dall'altra parte si ottiene la formula di \term{integrazione per parti}:
\[ \int_a^b f(x)\, g'(x) \dd x = \Big[ f(x)\, g(x) \Big]_a^b - \int_a^b f'(x)\, g(x) \dd x . \]
Serve quando l'integrando è un prodotto di due fattori, uno facile da derivare ($f$, che derivato si semplifica) e uno
facile da integrare ($g'$, di cui si conosce una primitiva $g$). Esempio:
\[ \int_0^1 x\, e^x \dd x = \Big[ x\, e^x \Big]_0^1 - \int_0^1 1 \cdot e^x \dd x = e - (e - 1) = 1 , \]
con $f = x$ (derivata 1) e $g' = e^x$ (primitiva $e^x$).
\end{richiamo}

\teorema[esponenzialemedia]{Media dell'esponenziale}{
  Se $X \sim \mathcal{E}(\lambda)$: \qquad $E[X] = \dfrac{1}{\lambda}$.
}

\begin{dimostrazione}
\passo{Si toglie il gradino} Il gradino annulla la parte negativa dell'asse:
\[ E[X] = \int_{-\infty}^{\infty} x\, f_X(x) \dd x = \int_{-\infty}^{\infty} x\, \lambda e^{-\lambda x} u(x) \dd x =
   \int_0^{\infty} x\, \lambda e^{-\lambda x} \dd x . \]
\passo{Per parti} Con $f(x) = x$ (derivata 1) e $g'(x) = \lambda e^{-\lambda x}$ (primitiva $g(x) = -e^{-\lambda x}$):
\[ \int_0^{\infty} x\, \lambda e^{-\lambda x} \dd x = \Big[ x \left( -e^{-\lambda x} \right) \Big]_0^{\infty} -
   \int_0^{\infty} \left( -e^{-\lambda x} \right) \dd x . \]
\passo{Il termine di bordo è nullo} In $x = 0$ vale $0$; per $x \to \infty$, $x\, e^{-\lambda x} \to 0$ (l'esponenziale
batte la potenza).
\passo{L'integrale rimasto}
\[ E[X] = \int_0^{\infty} e^{-\lambda x} \dd x = \left[ \frac{e^{-\lambda x}}{-\lambda} \right]_0^{\infty} = 0 -
   \frac{1}{-\lambda} = \frac{1}{\lambda} . \]
\end{dimostrazione}

\begin{figure}[H]
\centering
\begin{tikzpicture}
\begin{groupplot}[group style={group size=2 by 1, horizontal sep=1.6cm},
  quaderno, width=6.8cm, height=4.6cm, xmin=-0.6, xmax=4.4, xtick={0,1,2,3,4}]
\nextgroupplot[ylabel=$f_X(x)$, xlabel=$x$, ymin=0, ymax=2.3, ytick={0.5,1,2},
  legend style={font=\scriptsize, at={(0.97,0.97)}, anchor=north east}]
  \addplot[mspiA, domain=0:4.2, samples=80] {2*exp(-2*x)};
  \addplot[mspiB, domain=0:4.2, samples=80] {exp(-x)};
  \addplot[mspiC, domain=0:4.2, samples=80] {0.5*exp(-0.5*x)};
  \addplot[mspiA, line width=1.1pt, forget plot] coordinates {(-0.5,0) (0,0)};
  \legend{$\lambda = 2$, $\lambda = 1$, $\lambda = \frac12$}
\nextgroupplot[ylabel=$F_X(x)$, xlabel=$x$, ymin=0, ymax=1.2, ytick={0.5,1}, yticklabels={$\frac12$,$1$}]
  \addplot[mspiA, domain=0:4.2, samples=80] {1-exp(-2*x)};
  \addplot[mspiB, domain=0:4.2, samples=80] {1-exp(-x)};
  \addplot[mspiC, domain=0:4.2, samples=80] {1-exp(-0.5*x)};
  \addplot[mspigray, line width=1.1pt] coordinates {(-0.5,0) (0,0)};
\end{groupplot}
\end{tikzpicture}
\caption{Densità e CDF esponenziali. Con $\lambda$ grande la densità parte alta e scende in fretta (attese brevi,
media $\frac1\lambda$ piccola); con $\lambda$ piccolo è bassa e lunga (attese lunghe).}
\end{figure}

\begin{esempio}{L'esponenziale non ha memoria}
Il tempo di vita di un componente elettronico è $X \sim \mathcal{E}(\lambda)$. Il componente funziona da $s$ ore:
qual è la probabilità che funzioni per almeno altre $t$ ore? Con la probabilità condizionata e la CCDF:
\[
  P(X > s + t \mid X > s) = \frac{P(\{X > s + t\} \cap \{X > s\})}{P(X > s)} = \frac{P(X > s + t)}{P(X > s)} =
  \frac{e^{-\lambda(s + t)}}{e^{-\lambda s}} = e^{-\lambda t} = P(X > t) .
\]
(L'intersezione è $\{X > s + t\}$ perché chi supera $s + t$ supera anche $s$.) Il risultato non dipende da $s$: il
componente «non ricorda» quanto ha già lavorato, e la sua vita residua ha la stessa distribuzione di quella di un
componente nuovo. È la proprietà di \term{assenza di memoria}, che rende l'esponenziale il modello naturale dei
guasti casuali (non dovuti all'usura) e degli arrivi «a caso» di pacchetti o chiamate.
\end{esempio}

% ============================================================
\section{La variabile di Laplace}
% ============================================================

\dfn{Variabile di Laplace}{
  Una variabile aleatoria $X$ si dice \term{di Laplace} di parametro $\lambda > 0$, in breve $X \sim \mathcal{L}(\lambda)$,
  se ha densità
  \[ f_X(x) = \frac{\lambda}{2}\, e^{-\lambda \abs{x}} . \]
}

In questo caso $\mathrm{supp}[f_X(x)] = \R$: la variabile può assumere qualunque valore reale. La densità è fatta da due
esponenziali «schiena contro schiena», ciascuno dimezzato (così l'area totale è 1): è simmetrica attorno a 0, con una
punta nell'origine.

\teorema[laplacecdf]{CDF della variabile di Laplace}{
  Se $X \sim \mathcal{L}(\lambda)$:
  \[ F_X(x) = \begin{cases} \dfrac12\, e^{\lambda x} & x \leq 0 \\[2mm] 1 - \dfrac12\, e^{-\lambda x} & x \geq 0
     \end{cases} \]
}

\begin{dimostrazione}
\passo{Il valore assoluto} $\abs{t} = -t$ per $t \leq 0$ e $\abs{t} = t$ per $t \geq 0$: a sinistra dello zero la
densità è $\frac{\lambda}{2} e^{\lambda t}$, a destra $\frac{\lambda}{2} e^{-\lambda t}$.
\passo{Per $x \leq 0$} Si integra solo la parte sinistra. La primitiva di $e^{\lambda t}$ è $\frac{e^{\lambda t}}{\lambda}$,
che tende a 0 per $t \to -\infty$:
\[ F_X(x) = \int_{-\infty}^{x} \frac{\lambda}{2}\, e^{\lambda t} \dd t = \frac{\lambda}{2} \left[ \frac{e^{\lambda
   t}}{\lambda} \right]_{-\infty}^{x} = \frac12\, e^{\lambda x} . \]
\passo{Per $x \geq 0$} Si spezza l'integrale in zero, perché le due metà hanno formule diverse:
\[
  F_X(x) = \int_{-\infty}^{0} \frac{\lambda}{2}\, e^{\lambda t} \dd t + \int_0^x \frac{\lambda}{2}\, e^{-\lambda t} \dd t
  = \frac{\lambda}{2} \left( \left[ \frac{e^{\lambda t}}{\lambda} \right]_{-\infty}^{0} + \left[ \frac{e^{-\lambda
  t}}{-\lambda} \right]_0^x \right)
  = \frac{\lambda}{2} \left( \frac{1}{\lambda} + \frac{1 - e^{-\lambda x}}{\lambda} \right) .
\]
Raccogliendo: $\frac{\lambda}{2} \cdot \frac{2 - e^{-\lambda x}}{\lambda} = 1 - \frac12\, e^{-\lambda x}$. Si noti che le
due formule danno lo stesso valore, $\frac12$, in $x = 0$: metà della probabilità sta a sinistra dello zero e metà a
destra.
\end{dimostrazione}

Per la media non serve nessun calcolo, grazie a un'osservazione sulla simmetria.

\begin{richiamo}{funzioni pari e dispari}
Una funzione è \term{pari} se $g(-x) = g(x)$: il grafico è simmetrico rispetto all'asse verticale (esempi: $x^2$,
$\abs{x}$, $e^{-\abs{x}}$). È \term{dispari} se $g(-x) = -g(x)$: il grafico è simmetrico rispetto all'origine (esempi:
$x$, $x^3$, $\sin x$). Il prodotto di una funzione dispari per una pari è dispari: in particolare, se $f$ è pari,
$x\,f(x)$ è dispari.
\begin{center}
\begin{tikzpicture}
\begin{axis}[quaderno, width=7.4cm, height=4.2cm, xmin=-3.2, xmax=3.2, ymin=-0.65, ymax=0.65, xtick={-2,2},
  xticklabels={$-H$,$H$}, ytick=\empty, xlabel=$x$, grid=none]
  \addplot[fill=mspiA!25, draw=none, domain=0:2, samples=40] {x*exp(-abs(x))} \closedcycle;
  \addplot[fill=mspiB!25, draw=none, domain=-2:0, samples=40] {x*exp(-abs(x))} \closedcycle;
  \addplot[mspiric, domain=-3:3, samples=100] {x*exp(-abs(x))};
  \node[font=\scriptsize] at (axis cs:1,0.15) {$+$};
  \node[font=\scriptsize] at (axis cs:-1,-0.15) {$-$};
\end{axis}
\end{tikzpicture}
\end{center}
L'integrale di una funzione dispari su un intervallo simmetrico $[-H, H]$ è nullo: l'area a destra (positiva) e quella
a sinistra (negativa) sono uguali e opposte. Su tutta la retta, $\int_{-\infty}^{\infty}$, vale lo stesso
\textbf{purché} le due metà $\int_0^{\infty}$ e $\int_{-\infty}^{0}$ siano finite: se sono infinite, «$+\infty - \infty$»
non ha senso, e l'integrale non esiste.
\end{richiamo}

\teorema[laplacemedia]{Media della variabile di Laplace}{
  Se $X \sim \mathcal{L}(\lambda)$: \qquad $E[X] = \displaystyle\int_{-\infty}^{\infty} x\, \frac{\lambda}{2}\, e^{-\lambda
  \abs{x}} \dd x = 0$.
}

\begin{dimostrazione}
\passo{L'integrando è dispari} La densità è pari, quindi $x\, f_X(x)$ è dispari.
\passo{Le due metà sono finite} $\int_0^{\infty} x\, \frac{\lambda}{2} e^{-\lambda x} \dd x = \frac12 \cdot \frac1\lambda$,
perché è metà dell'integrale calcolato per la media dell'esponenziale. Allo stesso modo la metà sinistra vale
$-\frac{1}{2\lambda}$. La somma è zero.
\end{dimostrazione}

\begin{figure}[H]
\centering
\begin{tikzpicture}
\begin{groupplot}[group style={group size=2 by 1, horizontal sep=1.6cm},
  quaderno, width=6.8cm, height=4.4cm, xmin=-3.6, xmax=3.6, xtick={-3,-2,-1,1,2,3}]
\nextgroupplot[ylabel=$f_X(x)$, xlabel=$x$, ymin=0, ymax=0.62, ytick={0.5}, yticklabels={$\frac{\lambda}{2}$}]
  \addplot[mspiA, domain=-3.4:3.4, samples=121] {0.5*exp(-abs(x))};
\nextgroupplot[ylabel=$F_X(x)$, xlabel=$x$, ymin=0, ymax=1.2, ytick={0.5,1}, yticklabels={$\frac12$,$1$}]
  \addplot[mspiB, domain=-3.4:3.4, samples=121] {x<0 ? 0.5*exp(x) : 1-0.5*exp(-x)};
\end{groupplot}
\end{tikzpicture}
\caption{La variabile di Laplace con $\lambda = 1$: due esponenziali dimezzati, simmetrici attorno a 0.}
\end{figure}

% ============================================================
\section{La variabile di Cauchy}
% ============================================================

\dfn{Variabile di Cauchy}{
  Una variabile aleatoria $X$ si dice \term{di Cauchy} di parametri $(a, b)$, con $b > 0$, in breve
  $X \sim \mathcal{C}(a, b)$, se ha densità
  \[ f_X(x) = \frac{1}{b\pi}\; \frac{1}{1 + \left( \dfrac{x - a}{b} \right)^2} . \]
}

Anche qui $\mathrm{supp}[f_X(x)] = \R$. La densità è una campana simmetrica attorno ad $a$, che è quindi il
\emph{parametro di posizione}; $b$ ne regola la larghezza (a distanza $b$ dal centro la densità è scesa a metà del
massimo). Per calcolare la CDF servono due strumenti: il cambio di variabile e l'arcotangente.

\begin{richiamo}{la regola della catena e il cambio di variabile}
\textbf{La derivata di una funzione composta.} Se $y = g(x)$ e poi $z = f(y)$, la funzione composta è $z = f(g(x))$.
Quando $x$ cresce di un poco $h$, $y$ cresce di circa $g'(x)\,h$; e quando $y$ cresce di $g'(x)\,h$, $z$ cresce di
circa $f'(y) \cdot g'(x)\,h$. Le pendenze si moltiplicano:
\[ \frac{\dd}{\dd x} f\big(g(x)\big) = f'\big(g(x)\big)\; g'(x) . \]
È la \term{regola della catena}. Esempi: $\left( (1 + x^2)^3 \right)' = 3 (1 + x^2)^2 \cdot 2x$;
$\left( e^{-x^2/2} \right)' = e^{-x^2/2} \cdot (-x)$.

\textbf{Il cambio di variabile negli integrali.} Leggendo la regola della catena al contrario: se $G$ è una primitiva di
$f$, allora $G(g(x))$ ha derivata $f(g(x))\, g'(x)$, quindi
\[ \int_a^b f\big(g(x)\big)\, g'(x) \dd x = \int_{g(a)}^{g(b)} f(u) \dd u . \]
In pratica si pone $u = g(x)$, si sostituisce $\mathrm{d}u = g'(x) \dd x$ e si cambiano gli estremi. Il caso che ci
serve è quello lineare, $u = \frac{t - a}{b}$, con $\mathrm{d}u = \frac{\mathrm{d}t}{b}$, cioè $\mathrm{d}t = b \dd u$:
\[ \int_{-\infty}^{x} \phi\left( \frac{t - a}{b} \right) \dd t = b \int_{-\infty}^{\frac{x - a}{b}} \phi(u) \dd u . \]
Ha un significato geometrico semplice: spostare un grafico non cambia l'area, allargarlo di un fattore $b$ in
orizzontale la moltiplica per $b$.
\end{richiamo}

\begin{richiamo}{l'arcotangente}
La \term{tangente} di un angolo $\theta$ (in radianti) è $\tan\theta = \frac{\sin\theta}{\cos\theta}$. Sulla
circonferenza di raggio 1, è l'altezza a cui la semiretta inclinata di $\theta$ incontra la retta verticale $x = 1$.
Quando $\theta$ va da $-\frac{\pi}{2}$ a $\frac{\pi}{2}$ (da $-90^\circ$ a $90^\circ$, esclusi), la tangente assume una e una sola
volta ogni valore reale, da $-\infty$ a $+\infty$. Si può quindi invertire: l'\term{arcotangente}
$\arctan y$ è l'angolo, compreso fra $-\frac{\pi}{2}$ e $\frac{\pi}{2}$, la cui tangente vale $y$. Per esempio
$\arctan 0 = 0$, $\arctan 1 = \frac{\pi}{4}$, e $\arctan y \to \pm\frac{\pi}{2}$ per $y \to \pm\infty$.
\begin{center}
\begin{tikzpicture}[scale=1.5]
  \draw[-{Stealth}, mspigray] (-0.4,0) -- (1.6,0); \draw[-{Stealth}, mspigray] (0,-0.4) -- (0,1.9);
  \draw[mspigray] (1,0) arc (0:90:1);
  \draw[mspigray, dashed] (1,-0.3) -- (1,1.9) node[above, font=\scriptsize] {$x = 1$};
  \draw[mspiA!85!black, line width=0.8pt] (0,0) -- (1,1.03);
  \draw[mspiB!85!black, line width=0.8pt] (0,0) -- (1,1.38);
  \fill[evid, opacity=0.8] (0,0) -- (1,1.03) -- (1,1.38) -- cycle;
  \draw[mspiA!85!black, line width=0.8pt] (0,0) -- (1,1.03);
  \draw[mspiB!85!black, line width=0.8pt] (0,0) -- (1,1.38);
  \draw[mspiA!85!black] (0.35,0) arc (0:45.85:0.35); \node[font=\scriptsize] at (0.47,0.17) {$\theta$};
  \draw[line width=1.6pt, mspiric] (1,0) -- (1,1.03);
  \node[font=\scriptsize, mspiric, anchor=west] at (1.03,0.5) {$\tan\theta$};
  \node[font=\scriptsize, anchor=west] at (1.03,1.22) {$\mathrm{d}(\tan\theta)$};
  \node[font=\scriptsize, anchor=north west] at (0.66,0.6) {$r$};
  \fill (1,1.03) circle (0.025);
\end{tikzpicture}
\hspace{0.8cm}
\begin{tikzpicture}
\begin{axis}[quaderno, width=7cm, height=4.6cm, xmin=-6, xmax=6, ymin=-1.9, ymax=1.9, xtick={-4,4},
  ytick={-1.5708,1.5708}, yticklabels={$-\frac{\pi}{2}$,$\frac{\pi}{2}$}, xlabel=$y$, grid=none]
  \draw[mspigray, dashed] (axis cs:-6,1.5708) -- (axis cs:6,1.5708);
  \draw[mspigray, dashed] (axis cs:-6,-1.5708) -- (axis cs:6,-1.5708);
  \addplot[mspiric, domain=-6:6, samples=100] {rad(atan(x))};
  \node[font=\small, mspiric, anchor=north west] at (axis cs:0.6,0.9) {$\arctan y$};
\end{axis}
\end{tikzpicture}
\end{center}
\textbf{La sua derivata è $\frac{1}{1 + y^2}$.} Nel disegno a sinistra, la semiretta inclinata di $\theta$ incontra la
retta $x = 1$ a distanza $r = \frac{1}{\cos\theta}$ dall'origine, cioè $r^2 = 1 + \tan^2\theta$ (per Pitagora). Ruotando
la semiretta di un angolo piccolo $\mathrm{d}\theta$, il punto d'incontro sale di $\mathrm{d}(\tan\theta)$. L'arco
percorso a distanza $r$ è lungo $r\,\mathrm{d}\theta$ ed è perpendicolare alla semiretta; la retta verticale forma con
esso l'angolo $\theta$, quindi il tratto verticale è più lungo: $\mathrm{d}(\tan\theta) = \frac{r\,\mathrm{d}\theta}{\cos\theta}
= r^2\,\mathrm{d}\theta = (1 + \tan^2\theta)\,\mathrm{d}\theta$. Posto $y = \tan\theta$, cioè $\theta = \arctan y$:
\[ \frac{\mathrm{d}\theta}{\mathrm{d}y} = \frac{1}{1 + y^2} \qquad\Longrightarrow\qquad (\arctan y)' = \frac{1}{1 + y^2} . \]
Quindi $\arctan$ è una primitiva di $\frac{1}{1 + y^2}$, che è proprio la forma della densità di Cauchy.
\end{richiamo}

\teorema[cauchycdf]{CDF della variabile di Cauchy}{
  Se $X \sim \mathcal{C}(a, b)$: \qquad
  $F_X(x) = \dfrac12 + \dfrac{1}{\pi} \arctan\left( \dfrac{x - a}{b} \right)$.
}

\begin{dimostrazione}
\passo{Il cambio di variabile} Con $u = \frac{t - a}{b}$ si ha $\mathrm{d}t = b \dd u$, e l'estremo $t = x$ diventa
$u = \frac{x - a}{b}$:
\[ F_X(x) = \frac{1}{b\pi} \int_{-\infty}^{x} \frac{1}{1 + \left( \frac{t - a}{b} \right)^2} \dd t =
   \frac{1}{b\pi} \int_{-\infty}^{\frac{x - a}{b}} \frac{1}{1 + u^2}\, b \dd u = \frac{1}{\pi}
   \int_{-\infty}^{\frac{x - a}{b}} \frac{1}{1 + u^2} \dd u . \]
\passo{L'arcotangente} La primitiva è $\arctan u$, e $\arctan(-\infty) = -\frac{\pi}{2}$:
\[ F_X(x) = \frac{1}{\pi} \left[ \arctan\left( \frac{x - a}{b} \right) - \arctan(-\infty) \right] = \frac{1}{\pi}
   \left[ \arctan\left( \frac{x - a}{b} \right) + \frac{\pi}{2} \right] = \frac12 + \frac{1}{\pi} \arctan\left(
   \frac{x - a}{b} \right) . \]
\passo{Una verifica} Per $x \to +\infty$, $\arctan \to \frac{\pi}{2}$ e $F_X \to \frac12 + \frac12 = 1$: l'area totale è
1, quindi $f_X$ è davvero una densità.
\end{dimostrazione}

\begin{figure}[H]
\centering
\begin{tikzpicture}
\begin{groupplot}[group style={group size=2 by 1, horizontal sep=1.6cm},
  quaderno, width=6.8cm, height=4.4cm, xmin=-5.5, xmax=5.5, xtick={-4,-2,2,4}]
\nextgroupplot[ylabel=$f_X(x)$, xlabel=$x$, ymin=0, ymax=0.7, ytick={0.3183,0.6366}, yticklabels={$\frac1\pi$,$\frac2\pi$},
  legend style={font=\scriptsize, at={(0.98,0.98)}, anchor=north east}]
  \addplot[mspiA, domain=-5.3:5.3, samples=121] {1/(0.5*pi*(1+(x/0.5)^2))};
  \addplot[mspiB, domain=-5.3:5.3, samples=121] {1/(pi*(1+x^2))};
  \addplot[mspiC, domain=-5.3:5.3, samples=121] {1/(2*pi*(1+(x/2)^2))};
  \legend{$b = \frac12$, $b = 1$, $b = 2$}
\nextgroupplot[ylabel=$F_X(x)$, xlabel=$x$, ymin=0, ymax=1.2, ytick={0.5,1}, yticklabels={$\frac12$,$1$}]
  \addplot[mspiA, domain=-5.3:5.3, samples=121] {0.5+rad(atan(x/0.5))/pi};
  \addplot[mspiB, domain=-5.3:5.3, samples=121] {0.5+rad(atan(x))/pi};
  \addplot[mspiC, domain=-5.3:5.3, samples=121] {0.5+rad(atan(x/2))/pi};
\end{groupplot}
\end{tikzpicture}
\caption{Densità e CDF di Cauchy con $a = 0$. Le code scendono come $\frac{1}{x^2}$, molto più lentamente di un
esponenziale: valori lontanissimi dal centro sono rari, ma non abbastanza rari.}
\end{figure}

\subsection{La variabile di Cauchy non ha media}

La densità è simmetrica attorno ad $a$, e verrebbe da dire che la media è $a$. Invece la media \textbf{non è definita},
perché la funzione $x\, f_X(x)$ non è integrabile. Per vederlo serve una funzione che crescesse senza limite, ma molto
lentamente.

\begin{richiamo}{il logaritmo}
Il \term{logaritmo naturale} $\ln x$ (per $x > 0$) è l'inversa dell'esponenziale: $\ln x$ è il numero $y$ tale che
$e^y = x$. Quindi $\ln 1 = 0$, $\ln e = 1$, $\ln(e^{10}) = 10$, e trasforma i prodotti in somme:
$\ln(xy) = \ln x + \ln y$, $\ln\frac{x}{y} = \ln x - \ln y$ (perché $e^{a} e^{b} = e^{a + b}$). Cresce senza limite, ma
lentissimamente: per arrivare a 10 serve $x = e^{10} \approx 22\,026$.

\textbf{La sua derivata è $\frac1x$.} Se $y = \ln x$, allora $x = e^y$. Quando $y$ cresce di $\mathrm{d}y$, $x$ cresce di
$e^y\, \mathrm{d}y = x\, \mathrm{d}y$; quindi $\frac{\mathrm{d}y}{\mathrm{d}x} = \frac1x$. Ecco la primitiva che mancava
alla tabella del Capitolo~10, il caso $n = -1$: $\int_1^H \frac{1}{x} \dd x = \ln H$, che tende a infinito. L'area sotto
$\frac{1}{x}$ è infinita, mentre quella sotto $\frac{1}{x^2}$ è finita.

Con la regola della catena: $\frac{\dd}{\dd x} \ln(1 + x^2) = \frac{2x}{1 + x^2}$.
\end{richiamo}

\teorema[cauchymedia]{La variabile di Cauchy non ha media}{
  Se $X \sim \mathcal{C}(a, b)$, l'integrale $\int_{-\infty}^{\infty} x\, f_X(x) \dd x$ non esiste: le sue due metà
  valgono $+\infty$ e $-\infty$.
}

\begin{dimostrazione}[con $a = 0$ e $b = 1$]
\passo{La metà destra} Con la primitiva del richiamo:
\[ \int_0^H x\, \frac{1}{\pi}\, \frac{1}{1 + x^2} \dd x = \frac{1}{\pi} \left[ \frac12 \ln(1 + x^2) \right]_0^H =
   \frac{1}{2\pi} \ln(1 + H^2) \;\longrightarrow\; +\infty \qquad (H \to \infty) . \]
\passo{La metà sinistra} Per simmetria vale $-\infty$.
\passo{Conclusione} Le due metà non sono finite, e «$+\infty - \infty$» non è un numero. Il motivo è che per $x$ grande
$x\, f_X(x) \simeq \frac{1}{\pi x}$, che ha area infinita come $\frac1x$. Per $a$ e $b$ qualsiasi, il cambio di variabile
$u = \frac{x - a}{b}$ riporta a questo caso.
\end{dimostrazione}

Si può però definire un \textbf{punto di simmetria}, con un \term{integrale a valor principale}: invece di mandare a
infinito i due estremi in modo indipendente, li si allarga \emph{insieme}, simmetricamente.

\teorema[cauchyvp]{Il valor principale}{
  Se $X \sim \mathcal{C}(a, b)$:
  \[ \frac{1}{b\pi} \lim_{H \to \infty} \int_{-H}^{H} \frac{x}{1 + \left( \frac{x - a}{b} \right)^2} \dd x = a . \]
}

\begin{dimostrazione}
\passo{Il cambio di variabile} Con $u = \frac{x - a}{b}$, cioè $x = ub + a$ e $\mathrm{d}x = b \dd u$; gli estremi
$\mp H$ diventano $\frac{\mp H - a}{b}$:
\[ \frac{1}{b\pi} \lim_{H \to \infty} \int_{\frac{-H - a}{b}}^{\frac{H - a}{b}} \frac{ub + a}{1 + u^2}\, b \dd u =
   \frac{b}{\pi} \lim_{H \to \infty} \int_{\frac{-H - a}{b}}^{\frac{H - a}{b}} \frac{u}{1 + u^2} \dd u +
   \frac{a}{\pi} \lim_{H \to \infty} \int_{\frac{-H - a}{b}}^{\frac{H - a}{b}} \frac{1}{1 + u^2} \dd u . \]
\passo{Il primo limite è zero} Con la primitiva $\frac12 \ln(1 + u^2)$ e la proprietà dei logaritmi:
\[ \int_{\frac{-H - a}{b}}^{\frac{H - a}{b}} \frac{u}{1 + u^2} \dd u = \frac12 \ln\frac{1 + \left(\frac{H - a}{b}\right)^2}{1 +
   \left(\frac{H + a}{b}\right)^2} \;\longrightarrow\; \frac12 \ln 1 = 0 , \]
perché per $H$ grandissimo numeratore e denominatore della frazione sono quasi uguali (entrambi circa
$\frac{H^2}{b^2}$), e il loro rapporto tende a 1.
\passo{Il secondo limite} Gli estremi tendono a $-\infty$ e $+\infty$, e l'integrale diventa
$\int_{-\infty}^{\infty} \frac{1}{1 + u^2} \dd u = \left[ \arctan u \right]_{-\infty}^{\infty} = \pi$. Quindi il totale
è $0 + \frac{a}{\pi} \cdot \pi = a$.
\end{dimostrazione}

\info{Che cosa vuol dire, in pratica, «non ha media»}{
  Per una variabile con media, la media campionaria di tante osservazioni si stabilizza attorno alla media statistica
  (Capitolo~5). Per la Cauchy questo non succede: ogni tanto esce un valore così enorme da spostare da solo la media di
  migliaia di osservazioni. La figura confronta la media campionaria di osservazioni di Laplace e di Cauchy, entrambe
  simmetriche attorno a 0.
}

\begin{figure}[H]
\centering
\begin{tikzpicture}
\begin{axis}[width=12.5cm, height=5.2cm, xmode=log, xmin=1, xmax=10000, ymin=-2, ymax=6,
  xlabel={numero di osservazioni $n$ (scala logaritmica)}, ylabel={media campionaria},
  log ticks with fixed point, xtick={1,10,100,1000,10000}, tick label style={font=\scriptsize},
  label style={font=\small}, grid=major, grid style={mspigray!20}, axis lines=left,
  legend style={font=\scriptsize, at={(0.02,0.97)}, anchor=north west}]
  \draw[mspigray, dashed] (axis cs:1,0) -- (axis cs:10000,0);
  \addplot[mspiA, line width=1pt] coordinates {(1,0.109) (2,0.089) (3,-0.478) (4,-0.166) (5,-0.110) (6,-0.386) (7,-0.300) (8,-0.427) (10,-0.281) (12,-0.175) (15,-0.112) (18,0.090) (22,0.298) (26,0.189) (32,0.087) (38,0.090) (46,0.167) (56,0.071) (68,0.084) (83,0.149) (100,0.009) (121,-0.055) (147,-0.153) (178,-0.114) (215,-0.110) (261,-0.070) (316,-0.098) (383,-0.057) (464,-0.043) (562,-0.066) (681,-0.053) (825,-0.067) (1000,-0.055) (1212,-0.025) (1468,-0.038) (1778,-0.025) (2154,-0.044) (2610,-0.033) (3162,-0.022) (3831,-0.021) (4642,-0.005) (5623,-0.003) (6813,-0.006) (8254,-0.007) (10000,-0.006)};
  \addplot[mspiB, line width=1pt] coordinates {(1,-1.092) (2,-1.490) (3,0.271) (4,-0.095) (5,-0.408) (6,0.379) (7,0.522) (8,0.509) (10,0.720) (12,0.778) (15,0.999) (18,0.421) (22,5.586) (26,4.511) (32,3.962) (38,3.473) (46,2.638) (56,2.971) (68,1.124) (83,0.548) (100,0.491) (121,0.301) (147,0.526) (178,0.386) (215,0.053) (261,1.177) (316,0.499) (383,0.511) (464,0.786) (562,1.813) (681,2.306) (825,1.455) (1000,1.029) (1212,0.421) (1468,0.424) (1778,-0.256) (2154,3.456) (2610,2.874) (3162,2.295) (3831,4.544) (4642,3.059) (5623,2.439) (6813,2.019) (8254,1.599) (10000,1.363)};
  \legend{Laplace $\mathcal{L}(1)$, {Cauchy $\mathcal{C}(0, 1)$}}
\end{axis}
\end{tikzpicture}
\caption{Media campionaria di osservazioni simulate. Quella di Laplace converge a 0; quella di Cauchy continua a fare
salti improvvisi anche dopo migliaia di osservazioni, ognuno causato da un singolo valore enorme.}
\end{figure}

\begin{table}[H]
\centering
\renewcommand{\arraystretch}{1.7}
\begin{tabular}{@{}lllll@{}}
  \toprule
  \textbf{Variabile} & \textbf{Densità} & \textbf{CDF} & \textbf{Supporto} & \textbf{Media} \\
  \midrule
  uniforme $\mathcal{U}(a, b)$ & $\frac{1}{b - a}$ su $(a, b)$ & $\frac{x - a}{b - a}$ su $[a, b]$ & $[a, b]$ &
    $\frac{a + b}{2}$ \\
  esponenziale $\mathcal{E}(\lambda)$ & $\lambda e^{-\lambda x} u(x)$ & $(1 - e^{-\lambda x}) u(x)$ & $[0, \infty)$ &
    $\frac{1}{\lambda}$ \\
  Laplace $\mathcal{L}(\lambda)$ & $\frac{\lambda}{2} e^{-\lambda \abs{x}}$ & $\frac12 e^{\lambda x}$ ($x \leq 0$),
    $1 - \frac12 e^{-\lambda x}$ ($x \geq 0$) & $\R$ & 0 \\
  Cauchy $\mathcal{C}(a, b)$ & $\frac{1}{b\pi} \frac{1}{1 + ((x - a)/b)^2}$ & $\frac12 + \frac1\pi \arctan\frac{x -
    a}{b}$ & $\R$ & non esiste \\
  \bottomrule
\end{tabular}
\caption{Le variabili continue notevoli di questo capitolo (le varianze sono nel Capitolo~13).}
\end{table}

\gbox{In sintesi}{mspiA!80!black}{
\begin{itemize}
  \item \textbf{Uniforme} $\mathcal{U}(a, b)$: densità costante $\frac{1}{b-a}$, CDF lineare, media $\frac{a+b}{2}$.
  \item \textbf{Esponenziale} $\mathcal{E}(\lambda)$: $\lambda e^{-\lambda x} u(x)$, CDF $1 - e^{-\lambda x}$, media
        $\frac1\lambda$ (per parti); non ha memoria: è il modello dei tempi di attesa.
  \item \textbf{Laplace} $\mathcal{L}(\lambda)$: $\frac{\lambda}{2} e^{-\lambda\abs{x}}$, due esponenziali simmetrici;
        media 0 perché $x f_X(x)$ è dispari e le due metà sono finite.
  \item \textbf{Cauchy} $\mathcal{C}(a, b)$: CDF $\frac12 + \frac1\pi \arctan\frac{x-a}{b}$ (cambio di variabile e
        arcotangente). Non ha media: $x f_X(x) \sim \frac{1}{\pi x}$ e $\int \frac1x = \ln$ diverge; il valor principale
        vale $a$.
  \item Strumenti: $(e^{ax})' = a e^{ax}$, integrazione per parti, regola della catena e cambio di variabile,
        $(\arctan y)' = \frac{1}{1+y^2}$, $(\ln x)' = \frac1x$, integrali di funzioni dispari.
\end{itemize}
}
